%==============================================================================
%  Packages and notation.
%
%  Deliberately sober: Times text and maths, restrained headings, no colour, no
%  rules, no decoration.  The page should look like a journal proof, not like a
%  template.
%
%  Every symbol used more than twice in the manuscript has a macro here, so that
%  a change of notation is a one-line edit.
%==============================================================================

%------------------------------------------------------------- fonts, encoding --
% amsmath/amssymb must be loaded *before* newtxmath, or \Bbbk is defined twice.
\usepackage{amsmath,amssymb}
\usepackage[T1]{fontenc}
\usepackage[utf8]{inputenc}
\usepackage[english]{babel}
\usepackage{newtxtext}             % Times text
\usepackage{newtxmath}             % matching Times maths
\usepackage{microtype}

%------------------------------------------------------------------- page ------
\usepackage[a4paper,top=2.2cm,bottom=2.2cm,left=1.9cm,right=1.9cm]{geometry}
\setlength{\columnsep}{7mm}
\setlength{\parindent}{1.2em}      % indent, no inter-paragraph skip: journal style
\usepackage{flushend}              % balance the two columns on the last page

\usepackage[switch]{lineno}        % line numbers in the outer margin (two-column)
\usepackage{authblk}
\renewcommand\Authfont{\normalsize}
\renewcommand\Affilfont{\small\itshape}
\setlength{\affilsep}{0.6em}

%--------------------------------------------------------------- headings ------
% Small, bold, unnumbered-looking.  Nothing larger than \normalsize below the
% title.
\usepackage[explicit]{titlesec}
\titleformat{\section}{\normalsize\bfseries}{\thesection.}{0.6em}{#1}
\titleformat{\subsection}{\normalsize\itshape}{\thesubsection.}{0.6em}{#1}
\titleformat{\subsubsection}[runin]{\normalsize\itshape}{}{0em}{#1.}
\titlespacing*{\section}{0pt}{1.4ex plus .3ex}{0.7ex}
\titlespacing*{\subsection}{0pt}{1.2ex plus .3ex}{0.5ex}
\titlespacing*{\subsubsection}{0pt}{1.0ex plus .2ex}{0.5em}

%--------------------------------------------------------------------- math ----
% (amsmath and amssymb are loaded with the fonts above.)
\usepackage{bm}
\usepackage{mathtools}
\usepackage{siunitx}
\sisetup{range-units=single, per-mode=symbol}
% Two-column columns are narrow: keep display equations short, and prefer
% \begin{split} or multline over a single long line.
\allowdisplaybreaks

%------------------------------------------------------- floats and tables -----
\usepackage{graphicx}
\graphicspath{{figures/}}
\usepackage{booktabs}
\usepackage{multirow}
\usepackage{stfloats}              % allows [b] placement of figure*/table*
\usepackage[font=small,labelfont=bf,labelsep=period]{caption}
\usepackage{subcaption}
% Anything wider than a column goes in figure* / table* and floats to the top of
% a page.  Never scale a figure below ~80 % to force it into one column.

%----------------------------------------------------- references and links ----
\usepackage[round,authoryear]{natbib}
\bibpunct{(}{)}{;}{a}{}{,}
\setlength{\bibsep}{0pt plus 0.3ex}
\renewcommand{\bibfont}{\small}
\usepackage[hidelinks]{hyperref}   % no coloured boxes, no coloured text
\usepackage[capitalise,nameinlink]{cleveref}
\crefname{equation}{Eq.}{Eqs.}
\Crefname{equation}{Equation}{Equations}
\crefname{section}{Sec.}{Secs.}
\crefname{figure}{Fig.}{Figs.}
\crefname{table}{Table}{Tables}

%------------------------------------------------------------- draft markup ----
% \TODO{...}  placeholder --- set \draftfalse (or delete the text) at submission.
% Grey and bracketed rather than red: visible on screen, invisible in spirit.
\usepackage{xcolor}
\definecolor{todogrey}{gray}{0.42}
\newif\ifdraft\drafttrue
% \draftfalse hides every \TODO and the line numbers -- do that at submission.
\ifdraft
  \newcommand{\TODO}[1]{\textcolor{todogrey}{\textit{[#1]}}}
\else
  \newcommand{\TODO}[1]{}
\fi

%==============================================================================
%  Notation
%==============================================================================

%---------------------------------------------------------------- operators ----
\newcommand{\pd}[2]{\frac{\partial #1}{\partial #2}}
\newcommand{\pdd}[2]{\frac{\partial^{2} #1}{\partial #2^{2}}}
\newcommand{\dd}[2]{\frac{\mathrm{d} #1}{\mathrm{d} #2}}
\renewcommand{\d}{\mathrm{d}}
\newcommand{\ii}{\mathrm{i}}
\newcommand{\e}{\mathrm{e}}
\newcommand{\avg}[1]{\left\langle #1 \right\rangle}
\newcommand{\abs}[1]{\left\lvert #1 \right\rvert}
\newcommand{\norm}[1]{\left\lVert #1 \right\rVert}
\newcommand{\ip}[2]{\left\langle #1 , #2 \right\rangle}
\DeclareMathOperator*{\argmax}{arg\,max}
\DeclareMathOperator{\diag}{diag}
\DeclareMathOperator{\Real}{Re}
\DeclareMathOperator{\Imag}{Im}

%----------------------------------------------------------------- vectors -----
\newcommand{\vect}[1]{\bm{#1}}
\newcommand{\mat}[1]{\mathsf{#1}}
\newcommand{\uvec}{\vect{u}}
\newcommand{\xvec}{\vect{x}}
\newcommand{\nvec}{\vect{n}}

%--------------------------------------------------- flow / configuration ------
\newcommand{\Rey}{\mathit{Re}}
\newcommand{\Retau}{\Rey_{\tau}}
\newcommand{\Fro}{\mathit{Fr}}
\newcommand{\utau}{u_{\tau}}
\newcommand{\Ub}{U_{b}}
\newcommand{\Umax}{U_{\max}}
\newcommand{\Dr}{D_{r}}
\newcommand{\Rh}{R_{h}}
\newcommand{\yp}{y^{+}}
\newcommand{\zp}{z^{+}}

%----------------------------------------------------------- LES / averaging ---
\newcommand{\filt}[1]{\widetilde{#1}}
\newcommand{\test}[1]{\widehat{#1}}
\newcommand{\sgs}{\mathrm{sgs}}
\newcommand{\nut}{\nu_{t}}
\newcommand{\meanx}[1]{\overline{#1}}
\newcommand{\fluc}[1]{#1'}

%------------------------------------------------------------------- SPOD ------
\newcommand{\csd}{\mat{S}}
\newcommand{\Wmat}{\mat{W}}
\newcommand{\spodmode}{\vect{\psi}}
\newcommand{\spodev}{\lambda}
\newcommand{\nblk}{N_{\mathrm{b}}}
\newcommand{\nfft}{N_{\mathrm{f}}}
\newcommand{\kx}{k_{x}}
\newcommand{\qhat}{\hat{\vect{q}}}

%--------------------------------------------------- linear / resolvent --------
\newcommand{\Lop}{\mat{L}}
\newcommand{\Rop}{\mat{R}}
\newcommand{\Mop}{\mat{M}}
\newcommand{\Bop}{\mat{B}}
\newcommand{\Cop}{\mat{C}}
\newcommand{\Iop}{\mat{I}}
\newcommand{\fhat}{\hat{\vect{f}}}
\newcommand{\resu}{\vect{\phi}}
\newcommand{\resf}{\vect{\varphi}}
\newcommand{\gaink}{\sigma}

%--------------------------------------------------------------- shorthand -----
\newcommand{\ie}{i.e.\ }
\newcommand{\eg}{e.g.\ }
\newcommand{\cf}{cf.\ }
